% Terminal apparatus of the concise study: the scope note and the References
% for the sources this companion actually uses. The revision timestamp and the
% rights colophon follow from the shared generation record.

\section{Appendix: Scope and Qualifications}

{\footnotesize
\textbf{Edition, formulary and occurrence.} This concise study treats the Mass
\latin{Dominica decima septima post Pentecosten}, II classis, of the
\work{Missale Romanum}, editio typica, Typis Polyglottis Vaticanis 1962,
pp.~398--400, nos.~1602--1611, in the universal 1962 calendar and no particular
calendar. Its occurrence on 20~September 2026 rests on the Missal's own
calendar and general rubrics (RG 91, 111~b, 127~b; RGMR 427, 434~b, 494~b), and
agrees with the dated 2026 Ordos of the Fraternity of St Peter and of the
Institute of Christ the King in France, cited for the date only.

\textbf{Relation to the expansive study.} Every claim here is drawn from the
canonical study of this same formulary, \textit{The Seventeenth Sunday after
Pentecost}, which develops the three readings at length, treats each element in
its own setting and prints the appointed Latin with its public-domain English;
nothing evidence-dependent is added here. The four rows on the first page
orient a reader and do not replace those readings' own four senses; the map and
the dossier reprint the study's, with the dates, relation labels, alternatives
and unresolved states exactly as the generated chronology record gives them.

\textbf{Text, English and rights.} Every Latin form of the ten proper elements
was read for the study on the page images of the CMAA facsimile of the typical
edition and compared with the Pustet \work{Missale Romanum} of Ratisbon, 1862
and with the
Clementine Vulgate; the collation is in \texttt{propers/verified.md}. Two textual questions remain
open: the source of the Communion's \latin{omnes} before \latin{reges}, and the
Latin version behind the Offertory's compilation. The scriptural English is the
Douay--Rheims in Challoner's revision at the canonical verses, so that
it does not render the Missal's Latin wherever a chant or lesson adapts its
source; the English of the orations and of the Trinity Preface is the anonymous
translation printed by Cummiskey in Philadelphia in 1861. Neither is an
approved liturgical translation of the 1962 Missal; English renderings of
patristic Latin not from a named translation are glosses of the Latin beside
them. The Latin of the formulary rests on its public-domain antecedent, the
Pustet Missal (17 U.S.C. 103(b)); the Douay--Rheims, the 1861 English, the
Latin Fathers in Migne and the nineteenth-century translations are in the
public domain; the Marietti Latin of Aquinas on Matthew is quoted briefly from
a modern presentation; the \work{New American Bible} introduction and the
Ordos are cited, not reproduced.

\textbf{Reception and its limits.} The witnesses named here are those the
study's sweep read at the loci below. The Greek Fathers were read
in the Nicene and Post-Nicene and Ante-Nicene translations, with Chrysostom's
gloss on Eph 4:6 and Theodoret on Eph 4:1--6 also read in Migne's Greek;
Bellarmine only in O'Sullivan's abridged English of 1866, and that in a tracked
web transcription (ecatholic2000.com) and not the printing; Hilary on Ps~118 at
the strophes \latin{Ain} and \latin{Sade} alone. Some Latin was read in web
deliveries of the standard editions (augustinus.it, monumenta.ch, the Corpus
Thomisticum) or in optical text layers of printings (Rupert, Durandus,
Guéranger, Irenaeus, Trent), not on page images. Greek commentaries on the four
psalms, Greek exegesis of Daniel~9 and any Doctor on Daniel were not reached.
No witness is cited for a reading of the 1962 Mass as a whole: Rupert of Deutz
and William Durandus expounded an office of this Sunday whose Gospel was
Luke~14, and they are used for the chants their books share with it. The chant and sacramentary evidence is H.~A. Wilson's
editions and gregorien.info's view of Hesbert's \work{Antiphonale Missarum
Sextuplex}; no critical edition and no manuscript was consulted.

\textbf{Chronology and records.} The Date cells on the second page come from
the Scripture chronology corpus through the record generated for this formulary.
That record gives the four psalms only the modern critical boundary for the
Psalter, taken from the introduction to the Psalms in the \work{New American
Bible, Revised Edition} (summarized, not quoted), and holds no traditional date
for them; for Daniel, Matthew and Ephesians it holds only the
traditional dates reported by the \work{Catholic Encyclopedia}. Matthew's
critical comparison comes from a separate query of the same corpus; it does
not replace or merge the traditional alternatives. No critical date for
Daniel or Ephesians is supplied. The Gospel's narrated event
is undated. The bounded search, its loci, the disagreements and the negative
results are in \texttt{research/scope.md}, the reasoning behind the three
readings in \texttt{research/interpretations.md}, and this companion's
production record in \texttt{research/production-review.md}.
\par}

\section*{References}
\runninghead{References}

{\scriptsize
\subsubsection*{Liturgical books and Bibles}
\begin{itemize}\setlength{\itemsep}{0.15em}\setlength{\parskip}{0pt}
\item \work{Missale Romanum ex decreto Sacrosancti Concilii Tridentini restitutum, Summorum Pontificum cura recognitum}, editio typica (Typis Polyglottis Vaticanis, 1962), CMAA facsimile: \work{Rubricae generales} 91, 111, 127; \work{Rubricae generales Missalis} 427, 434, 494; pp.~293 (no.~1082), 397--400 (nos.~1594, 1602--1611), 410--416, 485--486 (nos.~2258, 2262), [57] (no.~4613).
\item \work{Missale Romanum} (Ratisbon: Pustet, 1862), pp.~341--342, \latin{Dominica XVII. post Pentecosten}.
\item \work{The Roman Missal, Translated into the English Language for the Use of the Laity} (Philadelphia: Eugene Cummiskey, 1861), pp.~431--433, xxviii--xxix.
\item The Holy Bible, Douay--Rheims version, revised by Bishop Richard Challoner: Deut 6:5; Ps 32; 75; 101; 109; 118; Dan 9; Mt 21--22; 26:1--2; Luke 14:9--11; Eph 1--4. \work{Biblia Sacra Vulgatae editionis} (Clementine Vulgate): Ps 32; 75; 101; 118; Dan 9; Mt 22; Eph 4.
\item H.~A. Wilson, ed., \work{The Gelasian Sacramentary} (Oxford, 1894), pp.~231--232 (Liber III, sect.~XIII) and 359 (appendix, \latin{Hebd.\ xxi post Pentecosten}), and \work{The Gregorian Sacramentary under Charles the Great}, HBS 49 (1915), pp.~174--175; R.-J. Hesbert, \work{Antiphonale Missarum Sextuplex} (1935), as indexed at gregorien.info, day 819.
\end{itemize}

\subsubsection*{Fathers, Doctors and saints}
\begin{itemize}\setlength{\itemsep}{0.15em}\setlength{\parskip}{0pt}
\item St Augustine, \work{Enarrationes in Psalmos}, Latin text of Migne, PL 36--37 (Wikisource and Corpus Corporum): on Ps~32, I, 12, and II, s.~2, 5, 15--18; Ps~101, s.~1, 2--3; Ps~109, 3--6. Latin on Ps~75, 16--18 and on Ps~118, s.~26, 5 and s.~28, 1 from augustinus.it.
\item St Augustine, \work{On Christian Doctrine} I.22.20--21, I.27.28, I.30.31--33 (Nicene and Post-Nicene Fathers, 1st ser., vol.~2); \work{Of the Morals of the Catholic Church} I.11.18 (1st ser., vol.~4); \work{On the Merits and Remission of Sins} II.13 (1st ser., vol.~5).
\item St Basil the Great, \work{On the Holy Spirit} 16.38 (Nicene and Post-Nicene Fathers, 2nd ser., vol.~8).
\item Cassiodorus, \work{Expositio psalmorum}, Latin (monumenta.ch): on Ps~32:6, 12; Ps~75:12--13; Ps~101:1--2; Ps~118:124, 137.
\item St John Chrysostom, \work{Homilies on Matthew} 71 (Nicene and Post-Nicene Fathers, 1st ser., vol.~10); \work{Homilies on Ephesians} 9--11 (1st ser., vol.~13), the gloss on Eph 4:6 also in Migne, PG 62, col.~80.
\item St Gregory the Great, \work{Homiliae in Evangelia} 38.10, Latin.
\item St Hilary of Poitiers, \work{Tractatus super Psalmos}, on Ps~118, \latin{Ain} 10 and \latin{Sade} 1, CSEL 22 (1891), pp.~500--501, 515.
\item St Irenaeus of Lyons, \work{Against Heresies} I.22.1, Ante-Nicene Fathers, vol.~1, p.~347.
\item St Jerome, \work{Commentarii in Matthaeum} IV, on 22:41--46 (PL 26, 166--167); \work{Commentarii in epistulam ad Ephesios} II, on 4:1--6 (PL 26, 492--496); \work{Commentarii in Danielem}, on 9:2--24 (PL 25, 678--683).
\item Theodoret of Cyrus, \work{Interpretatio epistolae ad Ephesios}, on 4:1--6 (Migne, PG 82, 532--533).
\item St Thomas Aquinas, \work{Super Evangelium S.~Matthaei lectura} c.~22, lect.~4, and \work{Postilla super Psalmos}, Ps~32, nn.~5, 11 (Corpus Thomisticum); \work{Super Ephesios} c.~4, lect.~1--2, and \work{Super Philippenses} c.~4, on 4:6 (Dessain, 1857), pp.~314--317, 398.
\item St Robert Bellarmine, \work{A Commentary on the Book of Psalms}, trans. J. O'Sullivan (1866, abridged), read in the eCatholic2000 transcription of that translation: on Pss~32:6, 12; 75:9--13; 101:2; 118:137.
\item Council of Trent, sess.~XIII, \work{Decretum de ss.~Eucharistia}, cap.~2 (Tauchnitz, 1887).
\end{itemize}

\subsubsection*{Liturgical commentators and chronology sources}
\begin{itemize}\setlength{\itemsep}{0.15em}\setlength{\parskip}{0pt}
\item Rupert of Deutz, \work{De divinis officiis} XII.17--18 (Migne, PL 170, 325--326).
\item William Durandus, \work{Rationale divinorum officiorum} VI.131 and the chapter on the Eighteenth Sunday after Pentecost (1568 printing).
\item Prosper Guéranger, \work{The Liturgical Year}, vol.~11, trans. Laurence Shepherd (James Duffy, 1900), pp.~372--387.
\item \work{The Catholic Encyclopedia} (New York, 1907--1912): F.~E. Gigot, ``Book of Daniel'' (1908); P. Ladeuze, ``Epistle to the Ephesians'' (1909); E. Jacquier, ``Gospel of St.~Matthew'' (1911); F. Prat, ``St.~Paul'' (1911); A. Durand, ``The New Testament'' (1912).
\item United States Conference of Catholic Bishops, \work{New American Bible, Revised Edition}, introductions to Psalms (superscriptions and chronological limit) and Matthew (authorship, date and place), official web edition inspected 21 September 2026; summarized, not reproduced.
\item Fraternité Saint-Pierre (France), \work{Ordo du mois}, and Institut du Christ Roi Souverain Prêtre (France), \work{Ordo}, entries for 20~September 2026.
\end{itemize}
\par}
